\subsection{Data Stream Model}
\label{ss:data_stream}

We consider a distributed system consisting of a server and a set of $n$ users, $U=\{u_1,\ldots,u_n\}$.
Each user $u_i\in U$ reports a value $v_{it}\in\Omega$ at each discrete timestamp
$t$, e.g., the user's location in a city.
We assume users' values are from a finite domain of size $d$, denoted by
$\Omega=\{\omega_1,\ldots, \omega_d\}$.
All users' reported values at time $t$ form a dataset $D_t$, which consists of
$n$ rows and $d$ columns with $D_t[i][j]=1$ if user $u_i$ reported value
$\omega_j$ at time $t$ and $D_t[i][j]=0$ otherwise.
We allow a user $u_i$ not to report her value at time $t$, i.e., $u_i$ is
inactive at $t$, and in this case, the $i$-th row of $D_t$ is zero.
Let $n_t$ denote the number of active users at time $t$.
Given $D_t$, the server can compute statistics of interest to publish.
For example, the server may continuously publish a vector
$\bc_t=[c_{1t},\ldots,c_{dt}]^\top$ at each time $t$, where $c_{jt}$ denotes the
number of users reporting the location $\omega_j\in\Omega$ at time $t$, and
publishing $\bc_t$ is useful for people to know the traffic congestion situation
in a city at time $t$.

However, the user's value may be private (e.g., her location), and publishing
statistics computed from this raw data can compromise privacy.
For example, suppose that there is only one active user in the system. In this case, the
published statistics directly reveal the user's private value.
In addition, the server is not assumed to be trusted (e.g., it is compromised
or contains software bugs), and it may leak users' private values.
Therefore, privacy-preserving techniques are required to protect users'
sensitive data from the server.
